\documentclass[a4paper,11pt,oneside]{report}

\usepackage[utf8]{inputenc}
\usepackage[T1]{fontenc}
\usepackage{listings}
\usepackage[french]{babel}
\usepackage[table]{xcolor}
\usepackage{graphicx}
\usepackage{geometry}
\geometry{a4paper, top=2.2cm, bottom=2.5cm, left=2.4cm, right=2.4cm, headsep=0.7cm}
\usepackage{amsmath,amssymb}
\usepackage{array}
\usepackage{caption}
\usepackage{subcaption}
\usepackage{lmodern}
\usepackage{booktabs}
\usepackage{microtype}
\usepackage{setspace}
\usepackage{enumitem}
\usepackage{titlesec}
\usepackage{fancyhdr}
\usepackage{needspace}
\usepackage[most]{tcolorbox}
\usepackage{pgfgantt}
\usepackage{pgf-umlsd}
\usepackage{float}
\usepackage[numbers]{natbib}
\usepackage{longtable}

\newenvironment{fquote}
  {\begin{quote}\itshape}
  {\end{quote}}

\newcommand{\mychapter}[2]{%
  \clearpage
  \chapter*{#2}%
  \markboth{#2}{#2}%
  \addcontentsline{toc}{chapter}{#2}%
  \vspace{-0.35em}
}

\usepackage{wrapfig}
\usepackage{tikz}
\usetikzlibrary{shapes.geometric,calc,arrows.meta,positioning,fit}
\usepackage{adjustbox}
\usepackage{xparse}
\usepackage{hyperref}
\hfuzz=3pt

\definecolor{Primary}{HTML}{1F3864}
\definecolor{Secondary}{HTML}{2E74B5}
\definecolor{Accent}{HTML}{ED7D31}
\definecolor{Alert}{HTML}{C00000}
\definecolor{SoftBg}{HTML}{F7FAFC}
\definecolor{TableHeader}{HTML}{2E74B5}
\definecolor{TableAlt}{HTML}{D9E2F3}
\definecolor{NoteBg}{HTML}{FFF2CC}
\definecolor{NoteBorder}{HTML}{BF8F00}

\newcommand{\tblhead}[1]{\textcolor{white}{\textbf{#1}}}

\hypersetup{
  colorlinks=true,
  linkcolor=Primary,
  citecolor=Primary,
  urlcolor=Primary,
  hypertexnames=false,
  pdftitle={Conception, implémentation et analyse critique de modèles de deep learning pour la détection de la dépression, la classification d'images et la traduction automatique},
  pdfauthor={KAFIF Imane}
}
\urlstyle{same}
\raggedbottom
\setstretch{1.08}
\setcounter{tocdepth}{1}
\setcounter{secnumdepth}{0}
\setlength{\parskip}{0.35em}
\setlength{\parindent}{1.1em}
\setlength{\headheight}{14pt}
\emergencystretch=1em

\setlist[itemize]{topsep=0.35em,itemsep=0.2em,leftmargin=1.6em}
\setlist[enumerate]{topsep=0.35em,itemsep=0.2em,leftmargin=1.8em}

\titleformat{name=\chapter,numberless}[display]
  {\normalfont\huge\bfseries\color{Primary}}
  {}
  {0pt}
  {}
  [\vspace{0.45em}\color{Primary}\titlerule]
\titleformat{\section}
  {\normalfont\Large\bfseries\color{Primary}}
  {}
  {0pt}
  {}
\titleformat{\subsection}
  {\normalfont\large\bfseries\color{Secondary}}
  {}
  {0pt}
  {}
\titlespacing*{\chapter}{0pt}{0pt}{1.0em}
\titlespacing*{\section}{0pt}{1.1em}{0.5em}
\titlespacing*{\subsection}{0pt}{0.9em}{0.35em}

\pagestyle{fancy}
\fancyhf{}
\fancyhead[L]{\nouppercase{\leftmark}}
\fancyhead[R]{\nouppercase{\rightmark}}
\fancyfoot[C]{\thepage}
\renewcommand{\headrulewidth}{0.4pt}
\renewcommand{\headrule}{\color{Primary}\hrule height\headrulewidth}
\renewcommand{\footrulewidth}{0pt}

\captionsetup{
  font=small,
  labelfont=bf,
  labelsep=period,
  skip=6pt
}

\newtcolorbox{block}[2][]{%
  enhanced,
  breakable,
  colback=Primary!4!white,
  colframe=Primary!50!black,
  colbacktitle=Primary!10!white,
  coltitle=black,
  fonttitle=\bfseries,
  sharp corners,
  boxrule=0.6pt,
  left=6pt,right=6pt,top=5pt,bottom=5pt,
  before skip=7pt,after skip=7pt,
  title={#2},
  #1
}
\newtcolorbox{alertblock}[2][]{%
  enhanced,
  breakable,
  colback=Alert!4!white,
  colframe=Alert!60!black,
  colbacktitle=Alert!10!white,
  coltitle=black,
  fonttitle=\bfseries,
  sharp corners,
  boxrule=0.6pt,
  left=6pt,right=6pt,top=5pt,bottom=5pt,
  before skip=7pt,after skip=7pt,
  title={#2},
  #1
}

\newcommand{\slidepurpose}[2]{%
\begin{tcolorbox}[
  enhanced,
  boxrule=0.4pt,
  sharp corners,
  colback=SoftBg,
  colframe=Primary!25,
  left=5pt,right=5pt,top=3pt,bottom=3pt,
  before skip=3pt,after skip=8pt
]
\footnotesize\textbf{Rôle :} #1 \hfill \textbf{Étape du projet :} #2
\end{tcolorbox}
}

\newcommand{\defitem}[2]{\textbf{#1} : #2\par\vspace{0.32em}}

\newcommand{\solutionfigure}[2]{%
\begin{center}
\includegraphics[width=0.96\linewidth,height=0.52\textheight,keepaspectratio]{#1}
\captionsetup{type=figure}
\captionof{figure}{#2}
\end{center}
}

\newcommand{\tikzsolutionfigure}[2]{%
\begin{center}
\begin{adjustbox}{max width=0.96\linewidth, max totalheight=0.52\textheight}
#1
\end{adjustbox}
\captionsetup{type=figure}
\captionof{figure}{#2}
\end{center}
}

\newif\ifreportfirstcolumn
\NewDocumentEnvironment{columns}{ O{} }{%
\par\medskip\noindent\reportfirstcolumntrue\ignorespaces
}{%
\par\medskip
}

\NewDocumentEnvironment{column}{ m }{%
  \ifreportfirstcolumn\reportfirstcolumnfalse\else\hfill\fi
  \begin{minipage}[t]{#1}\vspace{0pt}\raggedright
}{%
\end{minipage}\ignorespacesafterend
}

\RenewDocumentEnvironment{frame}{ O{} m }{%
\Needspace{8\baselineskip}%
\subsection{#2}
\noindent
}{%
\par\vspace{0.45cm}
}

\newcommand{\resultspath}{../results}
\newcommand{\safegraphic}[2][]{%
    \IfFileExists{#2}{\includegraphics[#1]{#2}}{%
        \fbox{\parbox[c][5cm][c]{0.85\linewidth}{\centering Figure non disponible.\\Relancer les expériences pour générer \texttt{#2}.}}%
    }%
}
\newcommand{\safetableinput}[1]{%
    \IfFileExists{#1}{\input{#1}}{%
        \textit{Tableau à remplacer après exécution du pipeline. Fichier attendu : \texttt{#1}.}%
    }%
}

\begin{document}

\begin{titlepage}
\centering
{\small Royaume du Maroc}\\
{\small HONORIS UNITED UNIVERSITIES}\par
\vspace{0.3cm}
\color{Primary}\rule{\linewidth}{1pt}\color{black}\par
\vspace{0.5cm}

\begin{minipage}{5cm}
\centering
\safegraphic[width=5cm]{logo.png}
\end{minipage}\hfill
\begin{minipage}{10cm}
\raggedleft
{\small \'Ecole Marocaine Des Sciences de L'Ing\'enieur}\\[0.1cm]
{\small D\'epartement Informatique}
\end{minipage}

\vspace{2cm}
{\large\bfseries Projet de fin de module}\\[0.3cm]
{\large Module : \textbf{Deep Learning}}\\[0.3cm]
{\large\bfseries Fili\`ere : 4IADO G2}\\[0.15cm]
{\normalsize P\'eriode : ann\'ee universitaire 2025--2026}

\vspace{1cm}
\color{Primary}\rule{\linewidth}{1pt}\color{black}\par
\vspace{0.6cm}
{\color{Primary}\LARGE\bfseries Conception, impl\'ementation et analyse\\critique de mod\`eles de deep learning\\pour la d\'etection de la d\'epression,\\la classification d'images et la traduction automatique}\par
\vspace{0.5cm}
{\large\itshape\color{Accent} MLP \;\textbullet\; CNN \;\textbullet\; RNN \;\textbullet\; LSTM \;\textbullet\; GRU \;\textbullet\; Seq2Seq}\par
\vspace{0.6cm}
\color{Primary}\rule{\linewidth}{1pt}\color{black}\par

\vspace{1.5cm}
\noindent
\begin{minipage}{0.5\textwidth}
\raggedright\large
\emph{R\'ealis\'e par :}\\
Mme KAFIF Imane
\end{minipage}
\begin{minipage}{0.4\textwidth}
\raggedleft\large
\emph{Ma\^itre p\'edagogique :}\\
Dr. HIDILA Zineb
\end{minipage}

\vfill
{\large EMSI Casablanca}
\end{titlepage}

\pagenumbering{roman}
\chapter*{Remerciements}
\addcontentsline{toc}{chapter}{Remerciements}
\thispagestyle{empty}

\vspace{1cm}

Je tiens à exprimer ma profonde gratitude à toutes les personnes qui ont contribué, de près ou de loin, à la réalisation de ce projet. \\[0.5cm]

Je remercie tout particulièrement \textbf{Dr. HIDILA Zineb} pour son encadrement, ses conseils avisés et son soutien tout au long de cette aventure.

Je souhaite également remercier \textbf{l'EMSI Casablanca} pour l'accompagnement pédagogique, l'environnement de travail et les ressources mises à disposition durant ce projet.

Enfin, une pensée sincère à ma famille et à mes proches pour leur soutien moral et leurs encouragements constants.

\vfill

\begin{flushright}
    \textit{KAFIF Imane}\\
    \textit{2025/2026}
\end{flushright}

\clearpage
\chapter*{Résumé}
Ce projet propose une étude expérimentale et critique de plusieurs architectures de deep learning appliquées à trois cas d'usage distincts. La première partie s'intéresse aux données tabulaires avec la détection de la dépression (Depression Professional Dataset) en comparant des modèles MLP selon leurs stratégies d'initialisation (Xavier, Gaussienne, Constante). La deuxième partie aborde la vision par ordinateur en implémentant des opérations de convolution manuellement, puis en construisant une architecture CNN de type LeNet-5 et une variante paramétrable pour l'étude d'ablation sur le jeu de données CIFAR-10. Enfin, la troisième partie explore la modélisation de séquences avec une comparaison structurelle (RNN, LSTM, GRU) et l'implémentation d'un système Seq2Seq pour la traduction automatique de l'anglais vers le français, évalué à l'aide de la perplexité et du score BLEU. Ce rapport met en évidence l'importance du choix architectural, des hyperparamètres et des méthodes de décodage (Greedy vs Beam Search) pour optimiser les performances des modèles selon la nature des données.

\chapter*{Mots-clés}
Deep learning, PyTorch, MLP, CNN, RNN, LSTM, GRU, Seq2Seq, détection de la dépression, CIFAR-10, traduction automatique, score BLEU, étude d'ablation.

\tableofcontents
\clearpage
\pagenumbering{arabic}

\chapter{Introduction générale}

\section{Contexte}
Le deep learning a profondément transformé la manière de modéliser des données complexes en apprenant automatiquement des représentations internes pertinentes \citep{goodfellow2016deep}. Cependant, cette efficacité apparente ne signifie pas qu’une architecture unique convienne à toutes les situations. Une table numérique de variables médicales, une image structurée sur une grille spatiale et une phrase ordonnée dans le temps n’imposent ni les mêmes biais inductifs, ni les mêmes contraintes de calcul, ni les mêmes mécanismes de généralisation.

\section{Problématique}
La question centrale de ce projet est donc la suivante : \emph{comment choisir et justifier une architecture de deep learning en fonction de la structure des données étudiées ?} Pour y répondre, trois familles de modèles ont été développées et comparées dans un même cadre expérimental :
\begin{itemize}[leftmargin=1.2cm]
    \item les MLP pour une tâche de classification tabulaire ;
    \item les CNN pour une tâche de classification d’images ;
    \item les architectures récurrentes puis un schéma encodeur--décodeur pour des séquences textuelles.
\end{itemize}

\section{Objectifs du projet}
Les objectifs sont doubles. Le premier est \textbf{ingénierie} : produire un code PyTorch propre, modulaire, reproductible et facile à exécuter. Le second est \textbf{scientifique} : comparer les performances, interpréter les résultats et discuter les limites réelles des modèles. Le projet vise donc autant la maîtrise de l’implémentation que la qualité de l’analyse critique.

\section{Méthodologie générale}
La méthodologie suivie est identique dans les trois parties :
\begin{enumerate}
    \item choix d’un jeu de données réel et suffisamment léger ;
    \item prétraitement rigoureux et séparation apprentissage / validation / test ;
    \item implémentation des architectures en PyTorch ;
    \item entraînement reproductible avec sauvegarde des résultats ;
    \item évaluation par des métriques adaptées ;
    \item visualisation et interprétation ;
    \item synthèse critique répondant à la question scientifique posée.
\end{enumerate}

\section{Organisation du rapport}
La Partie I traite les MLP sur données tabulaires. La Partie II traite les CNN sur images, avec calculs manuels et visualisation des cartes de caractéristiques. La Partie III traite la modélisation séquentielle avec RNN, LSTM, GRU et Seq2Seq. Enfin, une discussion transversale compare les architectures à travers la géométrie des données, avant de conclure sur les limites du travail et les prolongements possibles.


\mychapter{Partie I : Perceptron Multicouche (MLP)}{Partie I : Perceptron Multicouche (MLP)}

\section{Contexte et jeu de données}
\slidepurpose{Mettre en œuvre un MLP pour des données tabulaires et maîtriser l'inspection des poids.}{Partie 1 / Section 1}

Cette première partie se concentre sur l'analyse de données tabulaires pour la prédiction de la dépression à l'aide du \emph{Depression Professional Dataset}. Ce jeu de données, composé de 2054 observations et 11 variables, contient des informations sociodémographiques et professionnelles telles que l'âge, le genre, la pression au travail, la satisfaction professionnelle, et les habitudes de sommeil. L'objectif est de prédire la variable cible binaire \texttt{Depression} (Oui/Non).

Les données ont été nettoyées par suppression des valeurs manquantes, puis les variables catégorielles ont été transformées via un \texttt{LabelEncoder}. Le jeu de données a été scindé en un ensemble de test (20\,\%) et un ensemble d'entraînement/validation (80\,\%, lui-même divisé en 80/20). Enfin, une normalisation (\texttt{StandardScaler}) a été appliquée.

\section{Architectures MLP}

Deux approches de définition du perceptron multicouche ont été explorées.

\begin{columns}
\begin{column}{0.48\textwidth}
\begin{block}{Approche \texttt{nn.Sequential}}
Un empilement séquentiel direct :
\begin{itemize}
    \item \textbf{Entrée} : 10 caractéristiques
    \item \textbf{Couche cachée 1} : 64 neurones, ReLU
    \item \textbf{Couche cachée 2} : 32 neurones, ReLU
    \item \textbf{Sortie} : 2 neurones (classification)
\end{itemize}
\end{block}
\end{column}
\begin{column}{0.48\textwidth}
\begin{block}{Approche Orientée Objet (\texttt{nn.Module})}
Une classe personnalisée \texttt{DepressionMLP} permettant un contrôle plus fin lors du passage avant (\emph{forward pass}). L'architecture interne reste identique.
\end{block}
\end{column}
\end{columns}

L'inspection des poids via \texttt{state\_dict()} a permis de vérifier l'initialisation de la couche \texttt{fc1.weight} (moyenne : $-0.0062$, écart-type inhérent à l'initialisation par défaut de PyTorch).

\section{Impact de l'initialisation des poids}

Nous avons comparé trois stratégies d'initialisation pour évaluer leur impact sur la convergence du modèle (optimiseur Adam, taux d'apprentissage de 0.01, 50 époques).

\begin{center}
\renewcommand{\arraystretch}{1.3}
\begin{tabular}{l c}
\rowcolor{TableHeader} \tblhead{Stratégie d'initialisation} & \tblhead{Perte (Loss) finale (Époque 50)} \\
\rowcolor{TableAlt} \textbf{Xavier (Glorot)} & 0.0114 \\
\textbf{Gaussienne standard} ($\mu=0, \sigma=0.01$) & 0.0902 \\
\rowcolor{TableAlt} \textbf{Constante} (poids $= 1.0$) & 1.4322 \\
\end{tabular}
\end{center}

\begin{alertblock}{Observation critique}
L'initialisation constante échoue à briser la symétrie des neurones, empêchant l'apprentissage effectif. L'initialisation de Xavier s'avère la plus performante car elle maintient la variance des activations constante à travers les couches, particulièrement adaptée aux fonctions \texttt{ReLU}.
\end{alertblock}

\solutionfigure{figures/fig_p1_loss_init.png}{Comparaison de la convergence selon la stratégie d'initialisation}

\section{Évaluation des performances}

Le modèle optimal a été sauvegardé puis rechargé pour évaluation sur le jeu de test.

\begin{center}
\renewcommand{\arraystretch}{1.3}
\begin{tabular}{l c}
\rowcolor{TableHeader} \tblhead{Métrique} & \tblhead{Score} \\
\rowcolor{TableAlt} \textbf{Accuracy} (Précision globale) & 97.81\,\% \\
\textbf{Précision} (Vrais Positifs / Prédits Positifs) & 88.10\,\% \\
\rowcolor{TableAlt} \textbf{Recall} (Vrais Positifs / Total Positifs réels) & 90.24\,\% \\
\textbf{F1-Score} (Moyenne harmonique) & 89.16\,\% \\
\end{tabular}
\end{center}

\begin{columns}
\begin{column}{0.48\textwidth}
\solutionfigure{figures/fig_p1_confusion_matrix_2.png}{Matrice de confusion}
\end{column}
\begin{column}{0.48\textwidth}
\solutionfigure{figures/fig_p1_metrics_barplot.png}{Métriques de performance}
\end{column}
\end{columns}

\section{Discussion : Pertinence et limites du MLP}

\begin{block}{Question de Synthèse I : Pertinence et limites du MLP sur les données tabulaires}
\textbf{Analyse de l'étudiante :} Le Perceptron Multicouche (MLP) est l'architecture de référence pour les données tabulaires car ces dernières ne possèdent pas de relations topologiques ou spatiales a priori (contrairement aux images). La connectivité totale (\emph{fully connected}) permet de capter des combinaisons complexes de caractéristiques non linéaires à l'aide de fonctions d'activation appropriées.

Cependant, ses limites majeures incluent la sensibilité extrême au choix de l'échelle des variables (nécessité absolue d'une normalisation) et le risque élevé de surapprentissage (\emph{overfitting}) lorsque le nombre de variables et de couches augmente, en raison de l'explosion du nombre de paramètres apprenables.
\end{block}

\mychapter{Partie II : Réseaux Convolutifs (CNN)}{Partie II : Réseaux Convolutifs (CNN)}

\section{Contexte et jeu de données}
\slidepurpose{Exploiter l'indépendance spatiale des données visuelles avec les CNN.}{Partie 2 / Section 1}

La deuxième partie s'intéresse à la vision par ordinateur avec le jeu de données CIFAR-10. Contrairement aux données tabulaires de la première partie, les images possèdent une forte corrélation spatiale locale. CIFAR-10 contient 60 000 images couleur (RGB) de dimensions 32x32 réparties en 10 classes (avions, voitures, chats, etc.). Les données ont été prétraitées en tenseurs PyTorch et normalisées pour chaque canal.

\section{Compréhension des opérations fondamentales}

Afin de maîtriser les mécanismes internes des CNN, nous avons implémenté "from scratch" les opérations fondamentales de corrélation croisée 2D et de pooling (Max et Average) avant de les valider formellement contre l'API native de PyTorch.

\textbf{Implémentation manuelle de la corrélation 2D (\texttt{corr2d}) :}
\begin{verbatim}
def corr2d(X, K):
    h, w = K.shape
    Y = torch.zeros((X.shape[0] - h + 1, X.shape[1] - w + 1))
    for i in range(Y.shape[0]):
        for j in range(Y.shape[1]):
            Y[i, j] = (X[i:i+h, j:j+w] * K).sum()
    return Y
\end{verbatim}

Une comparaison directe avec \texttt{nn.Conv2d}, \texttt{nn.MaxPool2d} et \texttt{nn.AvgPool2d} via \texttt{torch.allclose} a permis de valider mathématiquement la conformité de nos boucles locales.

\section{Architecture LeNet5CIFAR}

Nous avons implémenté une architecture de référence inspirée du célèbre LeNet-5, modifiée pour admettre des images RGB et intégrant des améliorations modernes telles que la \texttt{BatchNorm2d} et le \texttt{Dropout}.

\begin{itemize}
    \item \textbf{Bloc 1 :} Conv2d (3$\rightarrow$16 filtres, 5x5, padding=2) $\rightarrow$ BatchNorm $\rightarrow$ ReLU $\rightarrow$ MaxPool2d(2x2)
    \item \textbf{Bloc 2 :} Conv2d (16$\rightarrow$32 filtres, 5x5, padding=0) $\rightarrow$ BatchNorm $\rightarrow$ ReLU $\rightarrow$ MaxPool2d(2x2)
    \item \textbf{Classifieur :} Flatten $\rightarrow$ LazyLinear(120) $\rightarrow$ ReLU $\rightarrow$ Dropout(0.3) $\rightarrow$ Linear(84) $\rightarrow$ Dropout(0.2) $\rightarrow$ Linear(10)
\end{itemize}

Après un entraînement sur 10 époques (optimiseur Adam, taux d'apprentissage de 0.001), la précision finale (\textbf{Accuracy}) sur les 10 000 images de test a atteint \textbf{72.93\,\%}. La perte finale s'est établie à 0.7681.

\solutionfigure{figures/fig_p2_train_curves.png}{Évolution de la Perte et de la Précision}

\section{Étude d'ablation avec \texttt{FlexibleCNN}}

Afin de quantifier l'impact de chaque choix de design, un modèle entièrement paramétrable (\texttt{FlexibleCNN}) a été conçu. Six scénarios ont été entraînés de façon systématique sur 3 époques pour comparer leur convergence initiale.

\begin{center}
\renewcommand{\arraystretch}{1.3}
\begin{tabular}{l c}
\rowcolor{TableHeader} \tblhead{Configuration Architecturale} & \tblhead{Précision Test (\%)} \\
\rowcolor{TableAlt} \textbf{Baseline (Standard)} ($p=2, s=1$, MaxPooling, 16 filtres) & 67.97\,\% \\
\textbf{Sans Padding} ($p=0$) & 67.05\,\% \\
\rowcolor{TableAlt} \textbf{Stride Fort} ($s=2$) & 63.54\,\% \\
\textbf{Average Pooling} & 66.86\,\% \\
\rowcolor{TableAlt} \textbf{Plus de Filtres} (64 filtres) & 71.21\,\% \\
\textbf{Avec Conv 1x1} (Compression des canaux) & 70.45\,\% \\
\end{tabular}
\end{center}

\begin{alertblock}{Analyse Architecturale}
Les résultats soulignent l'importance du MaxPooling pour préserver les gradients forts comparativement à l'Average Pooling. L'augmentation du nombre de filtres (capacité du réseau) apporte un net gain (71.21\%), tout comme l'introduction de convolutions 1x1 (Bottleneck structurel, 70.45\%), démontrant l'utilité des mélanges de canaux intra-pixels. Un stride fort en début de réseau provoque en revanche une perte d'information préjudiciable (baisse à 63.54\%).
\end{alertblock}

\section{Visualisation des Cartes de Caractéristiques (Feature Maps)}

Pour comprendre l'extraction hiérarchique, le modèle en mode évaluation a été appliqué à une image de la classe \emph{cat} (chat).

\solutionfigure{figures/fig_p2_cat_original.png}{Image originale dé-normalisée}

Les activations en sortie de la première couche convolutionnelle (16 filtres) et de la seconde couche (32 filtres) révèlent la transition des caractéristiques géométriques bas niveau (contours) vers des motifs plus abstraits.

\solutionfigure{figures/fig_p2_feature_maps_conv1.png}{Feature maps de la première couche de convolution (Conv1)}
\solutionfigure{figures/fig_p2_feature_maps_conv2.png}{Feature maps de la deuxième couche de convolution (Conv2)}

\mychapter{Partie III : Modélisation de Séquences}{Partie III : Modélisation de Séquences}

\section{Contexte et préparation textuelle}
\slidepurpose{Modéliser des données séquentielles avec dépendances temporelles.}{Partie 3 / Section 1}

La dernière partie s'attaque aux données textuelles (NLP) issues de critiques de type IMDb, pour une tâche de traduction automatique de l'anglais vers le français. Les séquences diffèrent fondamentalement des images car leur longueur est variable et l'information dépend de l'ordre temporel des mots. 

La préparation a nécessité la création de vocabulaires distincts (Source et Cible) basés sur la fréquence des mots, incluant des \textit{tokens} spéciaux (\texttt{<unk>}, \texttt{<pad>}, \texttt{<bos>}, \texttt{<eos>}), afin d'encoder les séquences textuelles en identifiants numériques.

\section{Comparaison structurelle : RNN, LSTM, GRU}

L'étude des états cachés a mis en évidence les différences d'architecture entre un simple RNN, une cellule LSTM et une cellule GRU. 

\begin{itemize}
    \item \textbf{RNN :} État simple \texttt{(1, 2, 32)}. Sujet au problème de disparition du gradient.
    \item \textbf{LSTM :} L'état est un tuple \texttt{(hn, cn)}, soit un \textit{hidden state} et un \textit{cell state} (chacun de taille \texttt{(1, 2, 32)}), agissant comme une "bande transporteuse" d'information longue distance.
    \item \textbf{GRU :} Version optimisée combinant les portes d'oubli et de mise à jour, avec un seul état caché \texttt{(1, 2, 32)}, offrant un compromis robuste entre performance et légèreté computationnelle.
\end{itemize}

\section{Architecture Seq2Seq et Entraînement}

Pour la traduction, un modèle encodeur-décodeur (\texttt{Seq2Seq}) basé sur des cellules GRU a été déployé. 

\textbf{Définition du processus de traduction (Extrait Seq2Seq) :}
\begin{verbatim}
class Seq2Seq(nn.Module):
    def forward(self, src, tgt):
        # 1. L'encodeur traite la phrase source en entier
        enc_outputs = self.encoder(src)
        
        # 2. Le decodeur s'initialise avec l'etat final de l'encodeur
        dec_state = self.decoder.init_state(enc_outputs)
        
        # 3. Le decodeur genere la sequence cible
        return self.decoder(tgt, dec_state, enc_outputs)
\end{verbatim}

L'entraînement sur 120 époques intègre le \textbf{Teacher Forcing} (alimentation de la vraie étiquette au temps $t-1$ pour générer la sortie au temps $t$) et le \textbf{Gradient Clipping} (\texttt{clip\_grad\_norm\_}) fixé à 1.0 pour prévenir l'explosion des gradients (problème structurel à la \textit{Backpropagation Through Time}). La perte moyenne finale (Loss) est descendue à 0.0037, garantissant une \textbf{perplexité parfaite de 1.00}.

\section{Stratégies de Décodage et Score BLEU}

Deux stratégies de génération ont été comparées pour la traduction d'une phrase de test : \textit{"this movie is great and fantastic"} dont la vérité terrain est \textit{"ce film est formidable et fantastique"}.

\begin{enumerate}
    \item \textbf{Décodage Glouton (Greedy Decoding)} : Le modèle sélectionne systématiquement le \textit{token} ayant la probabilité maximale au temps courant. Cette approche locale peut manquer l'optimum global.
    \item \textbf{Recherche par faisceau (Beam Search, $k=3$)} : Le modèle explore simultanément les $k$ séquences les plus probables, gardant une trace cumulative des probabilités pour favoriser une génération plus cohérente.
\end{enumerate}

Le \textbf{score BLEU} (Bilingual Evaluation Understudy), reposant sur une analyse des n-grammes avec pénalité de brièveté, a permis d'évaluer quantitativement les prédictions :

\begin{center}
\renewcommand{\arraystretch}{1.3}
\begin{tabular}{l c}
\rowcolor{TableHeader} \tblhead{Stratégie de décodage} & \tblhead{Score BLEU Obtenu} \\
\rowcolor{TableAlt} \textbf{Décodage Glouton} & 0.3333 \\
\textbf{Beam Search (k=3)} & 0.3333 \\
\end{tabular}
\end{center}

\begin{alertblock}{Analyse des prédictions}
Dans cette expérimentation, le score BLEU des deux méthodes est identique. Le décodage glouton a généré \textit{"ce film est formidable et fantastique et fantastique fantastique et"}, démontrant une tendance à la répétition (boucle infinie caractéristique des RNN). Le Beam Search, bien qu'il ait produit \textit{"ce film est formidable et et fantastique et fantastique fantastique"}, souffre du même défaut dans cette configuration, soulignant la complexité d'un arrêt prématuré \texttt{<eos>} correctement inféré.
\end{alertblock}

\solutionfigure{figures/fig_p3_recap.png}{Évaluation des stratégies de décodage (Exemples de métriques typiques NLP)}

\chapter{Discussion transversale finale}

\textbf{Comment le deep learning adapte-t-il ses architectures à la structure
des données --- tabulaire, image et séquentielle --- et pourquoi un même
paradigme d'apprentissage supervisé doit-il être décliné différemment selon la
géométrie, la dépendance locale, la temporalité et la représentation des
données ?}

\paragraph{Unité du paradigme.}
Les trois parties du projet partagent un même cadre : l'apprentissage
différentiable supervisé par minimisation d'une fonction de coût
$\mathcal{L}(\theta)$ via descente de gradient stochastique. Le graphe de
calcul, la rétropropagation et l'optimiseur restent identiques en substance.

\paragraph{Diversité des biais inductifs.}
Ce qui change d'une architecture à l'autre est le \emph{biais inductif} que
chacune encode :
\begin{itemize}
    \item le MLP suppose que les entrées sont un vecteur plat de
          $\mathbb{R}^d$, sans structure interne ;
    \item le CNN suppose une grille spatiale 2D avec invariance approximative
          par translation et hiérarchie locale ;
    \item le RNN/LSTM/GRU suppose un ordre séquentiel avec dépendances
          contextuelles causales.
\end{itemize}
Formellement, le MLP est invariant à toute permutation des dimensions d'entrée
(modulo réarrangement des poids), le CNN est équivariant aux translations
(sortie translatée si l'entrée est translatée), et le RNN est sensible à
l'ordre causal de la séquence.

\paragraph{Conséquences expérimentales observées.}
\begin{itemize}
    \item Sur les données tabulaires (11 variables, Depression Professional Dataset), le MLP
          atteint un F1 $> 0{,}89$ car les variables sont déjà des descripteurs
          informatifs.
    \item Sur les images ($32 \times 32$ RGB, CIFAR-10), le CNN 
          atteint près de 73\,\% d'accuracy car il exploite la localité
          spatiale et les hiérarchies de traits visuels.
    \item Sur les séquences textuelles (traductions type IMDb), les architectures récurrentes
          modélisent des dépendances temporelles que ni le MLP ni le CNN
          standard ne capturent.
\end{itemize}

\paragraph{Synthèse.}
Le deep learning n'est donc pas un bloc monolithique mais une famille de
méthodes dont la pertinence dépend de l'adéquation entre le biais inductif
du modèle et la géométrie des données. L'unité réside dans l'apprentissage
de représentations par optimisation différentiable ; la diversité réside dans
les contraintes architecturales choisies pour rendre cet apprentissage
structurellement efficient.

\chapter{Limites du travail}

Ce projet remplit les objectifs pédagogiques et expérimentaux du module, mais plusieurs limites doivent être explicitées.

\begin{itemize}[leftmargin=1.2cm]
    \item Les budgets d’entraînement ont été volontairement limités pour préserver l’exécutabilité sur ordinateur portable.
    \item La partie tabulaire ne compare pas le MLP à des modèles non neuronaux souvent très compétitifs sur ce type de données.
    \item La partie vision teste le réseau sur CIFAR-10, dataset utile pédagogiquement mais avec une résolution d'image ($32\times32$) très faible au regard de problèmes visuels modernes.
    \item La partie séquentielle repose sur un sous-corpus filtré, ce qui simplifie la tâche mais limite la couverture linguistique.
    \item Le système Seq2Seq n’intègre ni attention explicite, ni Transformer, ni embeddings pré-entraînés.
    \item Les résultats séquentiels doivent donc être interprétés comme une démonstration robuste de pipeline et de principes, non comme un système de traduction à l’état de l’art.
\end{itemize}

Ces limites n’invalident pas le projet ; elles en fixent simplement le périmètre scientifique et computationnel.


\chapter{Conclusion générale}

Ce projet a permis de construire un pipeline complet de deep learning couvrant trois régimes de données complémentaires : tabulaire, image et séquentiel. Sur les données tabulaires, le MLP s’est révélé efficace dès lors que les variables sont correctement normalisées et que l’architecture reste simple. Sur les images, les CNN ont permis d’étudier concrètement le rôle de la localité, du partage des poids, du padding, du stride et du pooling. Sur les séquences, les architectures récurrentes ont montré comment la mémoire contextuelle devient essentielle dès que l’ordre des observations conditionne la prédiction.

Au-delà des performances chiffrées, la principale conclusion est méthodologique : une bonne architecture n’est pas seulement un modèle puissant, c’est un modèle dont les hypothèses implicites sont compatibles avec la structure des données. Le deep learning apparaît alors non comme une recette universelle, mais comme une famille de méthodes dont la valeur dépend de l’ajustement entre représentation, géométrie et objectif d’apprentissage.

Le projet ouvre naturellement vers plusieurs prolongements : intégration de modèles tabulaires non neuronaux pour enrichir la comparaison, recours à des CNN plus profonds avec augmentation de données, ou encore ajout d’un mécanisme d’attention et d’un Transformer pour la traduction. Dans le cadre du module, l’ensemble obtenu constitue toutefois une base cohérente, rigoureuse et pleinement exploitable.



\appendix
\chapter{Annexes expérimentales}

\section{Commandes principales}
\begin{verbatim}
pip install -r requirements.txt
python run_notebooks.py
python run_notebooks.py --notebook 01_mlp_tabular.ipynb
python run_notebooks.py --notebook 02_cnn_images.ipynb
python run_notebooks.py --notebook 03_rnn_seq2seq.ipynb
\end{verbatim}
Les notebooks sont aussi exécutables individuellement depuis le dossier \texttt{notebooks/}.

\section{Organisation des sorties}
\begin{itemize}[leftmargin=1.2cm]
    \item \texttt{results/figures/mlp} : courbes, matrices de confusion et comparaison F1.
    \item \texttt{results/figures/cnn} : exemples d’images, courbes, matrices de confusion, cartes de caractéristiques.
    \item \texttt{results/figures/sequence} : courbes LM et Seq2Seq, comparaison de perplexité, effet du clipping.
    \item \texttt{results/tables} : export CSV et LaTeX des métriques.
    \item \texttt{results/saved\_models} : meilleurs poids sauvegardés.
\end{itemize}

\section{Choix d’implémentation}
Les choix suivants ont été appliqués dans tout le projet :
\begin{itemize}[leftmargin=1.2cm]
    \item graine aléatoire fixée ;
    \item détection automatique CPU/GPU ;
    \item création automatique des dossiers de sortie ;
    \item séparation stricte entraînement / validation / test ;
    \item sauvegarde des meilleurs modèles ;
    \item export des tableaux aux formats CSV et LaTeX.
\end{itemize}

\section{Remarque sur la reproductibilité}
Les résultats dépendent de l’environnement d’exécution, du budget de calcul disponible et des versions installées. Le dépôt inclut une configuration raisonnable pour obtenir des sorties reproductibles sur machine standard. En cas de relance avec d’autres hyperparamètres, les tableaux et figures du rapport peuvent être régénérés automatiquement.



\bibliographystyle{plainnat}
\bibliography{references}

\end{document}
